\DocumentMetadata{
  pdfversion=2.0,
  pdfstandard=ua-2,
  lang=en-US,
  tagging=on,
  tagging-setup={math/setup=mathml-SE,tabsorder=structure}
}
\documentclass[11pt,letterpaper]{article}
\usepackage[margin=0.68in,includefoot]{geometry}
\usepackage{newcomputermodern}
\usepackage{amsmath,amscd,mathtools}
\usepackage{tikz,adjustbox}
\providecommand{\square}{\mdlgwhtsquare}
\usepackage[hidelinks]{hyperref}
\hypersetup{
  pdftitle={Topology First-Year Exam, May 2015, Part 2},
  pdfauthor={University of Florida Department of Mathematics},
  pdfsubject={Graduate mathematics examination},
  pdfdisplaydoctitle=true,
  bookmarksopen=true
}
\setcounter{secnumdepth}{0}
\setlength{\parindent}{0pt}
\setlength{\parskip}{0.28em}
\setlength{\emergencystretch}{3em}
\setlength{\leftmargini}{2.15em}
\setlength{\leftmarginii}{2.65em}
\setlength{\leftmarginiii}{3em}
\pagestyle{empty}
\raggedbottom
\allowdisplaybreaks
\NewTaggingSocketPlug{block/list/label}{examproblem}{%
  \tagstructbegin{tag=Lbl}\tagstructend%
  \tagstructbegin{tag=\LBody}%
  \tagstructbegin{tag=\ProblemHeadingTag,title-o={\ProblemHeadingText}}%
  \tagmcbegin{tag=\ProblemHeadingTag,actualtext={\ProblemHeadingText}}%
  #2%
  \tagmcend%
  \tagstructend%
}
\newenvironment{examproblems}[2]{%
  \def\ProblemHeadingTag{#2}%
  \AssignTaggingSocketPlug{block/list/label}{examproblem}%
  \begin{enumerate}%
  \setcounter{enumi}{#1}%
  \renewcommand{\labelenumi}{\textbf{\theenumi.}}%
  \setlength{\topsep}{0.35em}%
  \setlength{\partopsep}{0pt}%
  \setlength{\itemsep}{0.48em}%
  \setlength{\parsep}{0.18em}%
}{\end{enumerate}}
\newenvironment{examsubparts}{%
  \AssignTaggingSocketPlug{block/list/label}{default}%
  \begin{enumerate}%
  \setlength{\topsep}{0.18em}%
  \setlength{\partopsep}{0pt}%
  \setlength{\itemsep}{0.16em}%
  \setlength{\parsep}{0.1em}%
}{\end{enumerate}}
\newenvironment{examinstructions}{%
  \par\begingroup\itshape
}{%
  \par\endgroup\vspace{0.18em}
}
\makeatletter
\renewcommand\subsection{\@startsection{subsection}{2}{0pt}%
  {-1.25ex plus -.25ex minus -.1ex}%
  {0.55ex plus .1ex}%
  {\normalfont\large\bfseries}}
\renewcommand{\@maketitle}{%
  \newpage\null\vskip -1em%
  {\footnotesize\bfseries
    \href{https://www.ufl.edu/}{University of Florida}, \href{https://math.ufl.edu/}{Department of Mathematics}\par}%
  \vskip -0.5em%
  \tagstructbegin{tag=H1,title={Topology First-Year Exam, May 2015, Part 2}}%
  \tagpdfparaOff%
  \tagmcbegin{tag=H1}%
    {\LARGE\bfseries\href{https://gma.math.ufl.edu/exams/topology-first-year/}{Topology First-Year Exam}, May 2015, Part 2\par}%
  \tagmcend%
  \tagpdfparaOn%
  \tagstructend%
  \vskip 0.25em%
  \tagmcbegin{artifact}%
    \hrule height 1pt%
  \tagmcend%
  \vskip 1em%
}
\makeatother
\title{Topology First-Year Exam, May 2015, Part 2}
\author{}
\date{}
\begin{document}
\maketitle
\thispagestyle{empty}
\begin{examinstructions}
Answer all questions and work all problems.
\end{examinstructions}
\begin{examproblems}{0}{H2}
\def\ProblemHeadingText{Problem 1.}
\item
Suppose that~\(X\) is a regular Lindelöf space. Show that~\(X\) is normal.
\def\ProblemHeadingText{Problem 2.}
\item
Show that there is a continuous \(f:[0,1]\to[0,1]^2\) such that~\(f\) is onto.
\def\ProblemHeadingText{Problem 3.}
\item
Suppose that \(f,g:X\to S^n\) are continuous maps such that for all \(x\in X\), \(f(x)\) and \(g(x)\) are not antipodal. Show that~\(f\) and~\(g\) are homotopic.
\def\ProblemHeadingText{Problem 4.}
\item
Show that the Sorgenfrey line~\(R\) is regular and Lindelöf.
\def\ProblemHeadingText{Problem 5.}
\item
Let \((X,x_0)\) be a pointed space. The trivial loop is the function \(f:[0,1]\to(X,x_0)\) such that \(f(t)\equiv x_0\) for all \(t\in[0,1]\). Let \(f:[0,1]\to(X,x_0)\) be a loop. Define \(f^{-1}\) by \(f^{-1}(t)=f(1-t)\). Show that \(f^{-1}\) is a loop. Show that \(f\circ f^{-1}\) is homotopic to the trivial loop.
\def\ProblemHeadingText{Problem 6.}
\item
Consider three loops \(f,g,h:[0,1]\to(X,x_0)\). Show that \((f\circ g)\circ h\) is homotopic to \(f\circ(g\circ h)\).
\def\ProblemHeadingText{Problem 7.}
\item
Let \((X,x_0)\) be a pointed space and let \(\pi_1(X,x_0)\) be the collection of homotopy classes of loops. Show that \(\pi_1(X,x_0)\) is a group.
\begin{examsubparts}
\item[\textnormal{(1)}]
Let \(\alpha=[a]\) and \(\beta=[b]\) be elements of \(\pi_1(X,x_0)\). Define \(\alpha\cdot\beta=[a\circ b]\). Show that \(\alpha\cdot\beta\) is well-defined.
\item[\textnormal{(2)}]
Show that the operation \(\alpha\cdot\beta\) is associative.
\item[\textnormal{(3)}]
Show that under the operation \(\alpha\cdot\beta\), the homotopy class of the trivial loop is the identity.
\item[\textnormal{(4)}]
Show that under the operation \(\alpha\cdot\beta\), the homotopy class of \(a^{-1}\) is the inverse of~\(\alpha\).
\end{examsubparts}
\def\ProblemHeadingText{Problem 8.}
\item
Show that \(\pi_1(S^1,1)=\mathbb{Z}\).
\def\ProblemHeadingText{Problem 9.}
\item
Show that \(\pi_1(S^n,1)=1\) for all \(n>1\).
\def\ProblemHeadingText{Problem 10.}
\item
Let \(n>1\). Let \(P^n=S^n/D\), where~\(D\) identifies~\(x\) with \(-x\) for all \(x\in S^n\). Show that \(\pi_1(P^n,x_0)=\mathbb{Z}_2\).
\def\ProblemHeadingText{Problem 11.}
\item
Consider the matrix
\[
M=\begin{bmatrix}1&2\\2&1\end{bmatrix}.
\]
 This represents a group homomorphism \(M:\mathbb{Z}^2\to\mathbb{Z}^2\). Show that there is a map \(f:\mathbb{T}^2\to\mathbb{T}^2\) such that \(f_*=M\), where \(f_*:\pi_1(\mathbb{T}^2)\to\pi_1(\mathbb{T}^2)\) is the homomorphism induced by~\(f\) on the fundamental group.
\def\ProblemHeadingText{Problem 12.}
\item
Let \(G=\langle a,b,c\mid a\cdot b^{-1}\cdot c\rangle\) be a finitely presented group with 3 generators and 1 relation.
\begin{examsubparts}
\item[\textnormal{(a)}]
Give a space~\(X\) whose fundamental group is the free group on three generators.
\item[\textnormal{(b)}]
Describe a space~\(Y\) whose fundamental group is the group~\(G\).
\end{examsubparts}
\def\ProblemHeadingText{Problem 13.}
\item
State the following theorems.

\par
Seifert–van Kampen Theorem.

\par
The Urysohn Metrization Theorem

\par
The Urysohn Lemma

\par
The Tietze Extension Theorem

\par
The Brouwer Fixed Point Theorem

\par
The Fundamental Theorem of Algebra

\par
The Jordan Curve Theorem

\par
The Arcwise Connectedness Theorem

\par
The Hahn–Mazurkiewicz Theorem
\end{examproblems}
\end{document}
